
%!TeX ts-program = pdflatex
\documentclass[]{papertexKS}
\usepackage[utf8]{inputenc}
\usepackage[spanish]{babel}
\usepackage{amsmath}
\usepackage{amssymb}
\usepackage{dirtytalk}
\usepackage{ks-biblatex}

\bibliography{string_theory_bib}

\begin{document}
\begin{news}{2}
    {Búsqueda de diferentes universos utilizando algoritmos de búsqueda metaeurísticos en la teoría K}
    {Brandon Marquez-Salazar, Angel Díaz-Pacheco, Cesar Damian}
    {Artículo en progreso}
    {teoría_cuerdas_y_metaheuristicos}

\nocite{*}

\noindent\textcolor{color}{\Large{Introduction}} \label{sec:intro}

La física es una ciencia cuyo propósito es el estudio de los diversos
fenómenos que ocurren en el universo, modelando matemáticamente su
comportamiento para comprender mejor cómo funcionan las mecánicas del
universo.

In quantum physics there are several approaches to study the behaviour of subatomic elements in an attempt to understand the fundamental interactions in the universe.

In quantum field theory, there are some problems that come up when trying to analyze dome phenomena at subatomic interaction level.
String theory is an approach intended to build mathematically feasible models of quantum \textit{\bfseries gravity}\cite{StringTheoryAndMTheory} by understanding fundamental particles as strings instead of puntual elements.

The interest in string theory is that from it emerges the relativist behaviour leading to the possibility of building the \textit{so-called} \textbf{theory of everything} which is, fast and loose, a theory which consistently permits to integrate the four fundamental forces which are \textit{\textbf{electromagnetic force}}, \textit{\textbf{strong atomic force}}, \textit{\textbf{weak atomic force}} and \textit{\textbf{gravity}}.

Regardless of the mathematical beauty of string theory, there are some downs that physicists face whenever a new approach intends to formalize a general string theory. An important one is that string theory, seen as a framework able to characterize an unimaginable amount of possible universes; this is called \textit{\textbf the landscape}.

Possible universes can be classified by it's geometry as \textit{\textbf{de Sitter}}, \textit{\textbf{flat}} or \textit{\textbf{anti de Sitter}}, each one with different physical structure and interactions. A \textbf{de Sitter} universe is one that expands at a constant acceleration, an anti \textbf{de Sitter} is an extremely symmetric universe where one can return to the start point going \textit{entlang}, and a \textbf{flat} universe is an universe which is not expanding.
Following Einstein's relativity our universe is anti de Sitter.

From the landscape and geometric universe classification, \textit{\textbf{swapland}} conjectures are formulated. The \textit{\textbf{swapland}} mainly indicates that string theory may not be able to construct stable \textit{\textbf{de Sitter}} universes.

Other of the interesting observed behaviour in our universe is the presence of two type of quantum particles called \textbf{\textit{bosons}} which \say{conform energy} and fermions which \say{conform matter}. For the current particle model, there are differences between bosons and fermions, which are not necessarily symmetric; understanding as symmetry the existence of an energy particle per matter particle and vice versa. This constrains the search even more, so the configuration should let an universe which expands and also, have the particles detected in our standard particle model. In other words, finding an universe configuration which lets the $\lambda$CDM model emerge from it's structure.

In a previous from Cesar Damian \cite{Damian2022} work there was made an attempt to check the validity the conjectures in k-theory model finding metastable vacua, which confirms such conjectures using simulated annealing and ML algorrithms. However, since there are many posible configurations, there could be more possible configurations, there are plenty of vacua to be mapped, and since simulated annealing was the only one metaheuristic algorithm considered. The main proposal is a more algorithm-focused research testing different methods in order to map the most possible universe configurations.


\noindent\textcolor{color}{\Large{Methods}} \label{sec:methods}

\noindent\textcolor{color}{\Large{Results}} \label{sec: results}

\noindent\textcolor{color}{\Large{Discussion}} \label{sec: discussion}

\noindent\textcolor{color}{\Large{Conclusions}} \label{sec: conclusions}

\printbibliography[]
\end{news}
\end{document}
